\chapter{V5 JEPA Architecture}
\label{ch:architecture}

This chapter specifies the model architecture used throughout the
present work, an adaptation of the V-JEPA-2-AC
template~\cite{vjepa2_2025,assran2024vjepa} to a single-channel,
signed pseudoscalar field of variable spatial extent and 12-minute
cadence.

\section{Joint-embedding predictive doctrine}
\label{sec:arch-doctrine}

In a Joint-Embedding Predictive Architecture (JEPA), prediction is
performed in the embedding space of a learned encoder, not in pixel
space~\cite{assran2023ijepa,lecun2022path}. Two encoders are
maintained: a \emph{context encoder} that consumes a masked view of
the input and a \emph{target encoder} that consumes the
corresponding unmasked view. A predictor maps context embeddings to
target embeddings; the smooth-$L_{1}$ loss between the two governs
training. The target encoder is updated by exponential moving
average (EMA) of the context encoder's parameters and receives no
gradient. The asymmetry between the two streams prevents the
representational collapse that would otherwise drive both encoders
to a constant solution; no contrastive negatives are required.

Compared with pixel-space reconstruction
(masked autoencoder, MAE~\cite{he2022mae}; VideoMAE~\cite{tong2022videomae}; the Pathak
context encoder~\cite{pathak2016contextencoders}), the
embedding-space objective discards by construction the high-frequency
content of the target field that the encoder has chosen not to keep.
This addresses the persistence collapse observed in the V4
pixel-space forecaster on the same corpus.

\begin{figure}[htbp]
\centering
\begin{tikzpicture}[
  >=stealth, node distance=0.9cm and 0.9cm,
  font=\small,
  block/.style={draw, rounded corners, align=center, minimum width=2.4cm,
                minimum height=0.8cm, fill=blue!5},
  enc/.style={block, fill=blue!12},
  pred/.style={block, fill=orange!18},
  tgt/.style={block, fill=green!15},
  loss/.style={block, fill=red!12, minimum width=3.4cm},
  io/.style={align=center, font=\small\itshape},
  arr/.style={->, thick},
  ema/.style={->, dashed, thick, blue!60!black}]

  \node[io] (xin) {input cube $x \in \mathbb{R}^{B \times T \times 1 \times H \times W}$};
  \node[block, below=0.5cm of xin] (mask) {masking $M$\\(\S\ref{sec:arch-mask})};
  \node[enc, below left=0.9cm and 1.4cm of mask] (ctx) {context encoder\\ViT-Small + RoPE-2D};
  \node[tgt, below right=0.9cm and 1.4cm of mask] (tgt) {target encoder\\EMA, \texttt{stop\_grad}};
  \node[pred, below=1.1cm of ctx] (pred) {block-causal\\predictor};
  \node[io, below=0.5cm of pred] (zhat) {$\hat z_{\text{masked}}$};
  \node[io, below=0.5cm of tgt] (ztgt)  {$z_{\text{tgt}}$};
  \node[loss, below right=1.5cm and -0.9cm of pred] (lo) {smooth-$L_1$ over\\
        $M_{\text{loss}} = M_{\text{mask}} \wedge M_{\text{valid}} \wedge M_{\text{pad}}$};

  \draw[arr] (xin) -- (mask);
  \draw[arr] (mask) -- node[midway, above left, font=\scriptsize] {visible} (ctx);
  \draw[arr] (mask) -- node[midway, above right, font=\scriptsize] {full} (tgt);
  \draw[arr] (ctx) -- (pred);
  \draw[arr] (pred) -- (zhat);
  \draw[arr] (tgt) -- (ztgt);
  \draw[arr] (zhat) -| (lo);
  \draw[arr] (ztgt) -| (lo);
  \draw[ema] (ctx.north east) to[bend left=18]
       node[midway, above, font=\scriptsize, blue!60!black]
       {EMA, $\theta_{\text{tgt}} \leftarrow \tau\,\theta_{\text{tgt}}+(1-\tau)\,\theta_{\text{ctx}}$}
       (tgt.north west);
\end{tikzpicture}
\caption{V5 joint-embedding predictive architecture. The input
cube is masked once and routed to both encoders: the context
encoder sees only the visible tokens, the target encoder sees the
full input under \texttt{stop\_grad}. The block-causal predictor
maps context embeddings to the masked target positions; the
smooth-$L_{1}$ loss is computed over the mask intersection
$M_{\text{loss}} = M_{\text{mask}} \wedge M_{\text{valid}}
\wedge M_{\text{pad}}$ (\S\ref{sec:arch-loss}). The dashed arrow
records the parameter EMA from context to target with decay
$\tau$ (no gradient flows backward through it).}
\label{fig:jepa-arch}
\end{figure}
\FloatBarrier

\section{Input adapter}
\label{sec:arch-adapter}

Active-region cubes arrive as tensors $x \in \mathbb{R}^{B \times T
\times 1 \times H \times W}$, with $T = t_{\text{in}} + t_{\text{out}}$
frames at 12-minute cadence and $(H, W)$ varying per cube. A
bucketed sampler groups cubes of similar spatial extent within each
mini-batch. The input adapter reflection-pads $(H, W)$ to the next
multiple of the patch size $P = 16$ and lifts the single physical
channel (winding flux) to 13 latent channels via a $1 \times 1$
convolution. The latent channel count matches the channel budget of
the published Surya backbone and was kept for compatibility with the
abandoned Path~A. A boolean token-pad mask $M_{\text{pad}} \in
\{0,1\}^{h_{p} \times w_{p}}$ is propagated alongside the embedded
tensor, with $h_{p} = H_{p}/P$ and $w_{p} = W_{p}/P$, and is later
intersected with the loss mask so that padded positions never enter
the gradient.

\section{Context encoder}
\label{sec:arch-ctx}

The context encoder is a Vision Transformer (ViT) of the
\textsc{ViT-Small} class~\cite{dosovitskiy2020vit}. Each frame is patch-embedded by a
$\text{Conv2d}(13 \to d, P{=}16, s{=}16)$ stem with $d = 384$,
then processed by $L_{\text{ctx}} = 12$ pre-norm transformer blocks
of bidirectional spatial attention. Weights are shared across the
frame axis: each of the $T$ frames in a window is encoded
independently, so that all temporal mixing is concentrated in the
predictor (\S\ref{sec:arch-pred}).

Spatial position is encoded by a two-dimensional rotary position
embedding (RoPE)~\cite{su2021rope} whose coordinates are expressed in physical units of
megametres on the photospheric surface. With a pixel scale of $0.364
\,\mathrm{Mm/px}$ and patch size $P$ the angular frequency table is
fixed across cubes, so that the same physical separation between two
patches receives the same rotation regardless of where it sits within
a variable-sized cube. The encoder receives gradients via the
standard backward pass.

\section{Target encoder}
\label{sec:arch-tgt}

The target encoder is a parameter-wise EMA copy of the context
encoder; its forward pass is executed under
\texttt{torch.no\_grad}, so that no gradient flows through the target
branch. After each optimiser step the parameter update
\begin{equation}
\theta_{\text{tgt}} \;\leftarrow\; \tau \,\theta_{\text{tgt}}
\;+\; (1-\tau)\,\theta_{\text{ctx}},
\qquad \tau \in (0,1),
\label{eq:ema}
\end{equation}
is applied. The decay coefficient $\tau$ is the
single most consequential hyperparameter of the pretraining stage and
is swept in Chapter~\ref{ch:experiments}. The default value
$\tau = 0.996$ is inherited from V-JEPA~2~\cite{vjepa2_2025}; an
ablation between $\tau = 0.990$ and $\tau = 0.9995$ is reported in
\S\ref{sec:exp-ema}. The target encoder consumes the
\emph{clean} (unmasked) latent tensor, so that target embeddings are
not contaminated by the mask token.

\section{Block-causal predictor}
\label{sec:arch-pred}

The predictor is a $L_{\text{pred}} = 6$ block transformer of
embedding dimension $d = 384$. It receives the context embedding
tensor reshaped to $[\,B,\, T\,h_{p}\,w_{p},\, d\,]$ and applies a
three-dimensional rotary position embedding indexed by the triplet
(time in minutes, $y$ in Mm, $x$ in Mm). Within a frame, attention
is fully bidirectional; across frames, the attention mask is
block-causal, that is, frame $t$ may attend to frames $0,1,\ldots,t$
but not to frames $t+1,\ldots,T-1$. The block-causal structure is
the ``\textsc{ac}'' component of V-JEPA-2-AC and is necessary for
the future-block mask of \S\ref{sec:arch-mask} to constitute a
non-trivial pretext task.

\section{Mask catalogue and curriculum}
\label{sec:arch-mask}

Five mask families are catalogued in
\texttt{solarflare\_data/mask\_catalog.py}, summarised in
Table~\ref{tab:masks}. The default training mix is a categorical
distribution over three of the five families, $\{\text{tube}:0.5,\
\text{future}:0.3,\ \text{cross\_time}:0.2\}$. The tail-only and
random-region families are reserved for the curriculum and the
ablation, respectively.

\begin{table}[htbp]
\centering\small
\begin{tabular}{l l l l}
\toprule
Family & Shape & Scale & Purpose \\
\midrule
Short tube       & spatial block, all $T$           & 15\% area        & local AR structure \\
Long tube        & spatial block, all $T$           & 40\% area        & global flux motion \\
Future block     & full spatial, last $K$ frames    & $100\% \cdot K/T$ & future prediction \\
Cross-time       & full spatial, random $T$ frames  & 30\% time        & non-causal infill \\
Tail (downstream)& full spatial, last $t_{\text{out}}$ & $t_{\text{out}}/T$ & deployment-aligned \\
\bottomrule
\end{tabular}
\caption{V5 mask catalogue. Appendix~\ref{app:masks} renders one
realisation of each family on a single frame.}
\label{tab:masks}
\end{table}

The curriculum interpolates the mix over training: from epoch $0$ to
$0.25\,E$ the model sees the \emph{tail-only} family ($\{\text{tail}:
1\}$); from $0.25\,E$ to $0.40\,E$ the mix interpolates linearly
toward the default; from $0.40\,E$ to the end of training the
default mix is used. The slow schedule
($0.25 / 0.40$) was selected over a faster schedule ($0.10 /
0.20$) on the basis of finding~F1 (see
Chapter~\ref{ch:experiments}). The intent of the curriculum is to
let the encoder consolidate the deployment-aligned tail mask before
the harder tube and future families are introduced.

\section{Loss}
\label{sec:arch-loss}

The pretraining objective is the smooth-$L_{1}$ (Huber)
distance~\cite{girshick2015fastrcnn} between the predicted and
target embeddings, evaluated only at token positions
that satisfy three constraints simultaneously: the position is masked
in the context view, its underlying pixels are physically valid
(after the winding-flux outlier guard of Chapter~\ref{ch:data}),
and the position is not a padded token from
\S\ref{sec:arch-adapter}. The intersection
\begin{equation}
\mathcal{M}_{\text{loss}} \;=\; M_{\text{mask}}
\,\wedge\, M_{\text{valid}}
\,\wedge\, M_{\text{pad}}
\label{eq:lossmask}
\end{equation}
is computed token-wise and the loss is averaged over $|\mathcal{M}_{\text{loss}}|$.
A subtle bug in the token-pad mask used in earlier sanity runs ---
the trailing-frame slice $[\![-0,]\!]$ was a Python no-op and
disabled the pad mask entirely --- was fixed prior to the experiment
campaign of Chapter~\ref{ch:experiments}; the fix has no effect on
runs in which no padding was applied.

\section{Parameter budget}
\label{sec:arch-scale}

\begin{table}[htbp]
\centering\small
\begin{tabular}{l r c}
\toprule
Component & Parameters & Gradient \\
\midrule
Input adapter            & \phantom{0}0.2~M & yes \\
Context encoder (ViT-S, $L{=}12$, $d{=}384$) & 22~M  & yes \\
Target encoder (EMA copy of context) & 22~M  & no  \\
Block-causal predictor ($L{=}6$, $d{=}384$)  & \phantom{0}9~M  & yes \\
\midrule
Total                                         & 55~M  &     \\
Trainable                                     & 33~M  &     \\
\bottomrule
\end{tabular}
\caption{V5 parameter budget. Trainable count excludes the EMA target.}
\label{tab:scale}
\end{table}

\noindent The full forward graph in PyTorch fits in
approximately $12\,\mathrm{GB}$ of CUDA memory at batch~$\times$~time
$= 4 \times 14$, single precision activations under
\texttt{bf16} autocast and gradient checkpointing on the context
encoder.
